\subsection{Self-adjoint operators defined by a positive partial Hermitian form}

In this issue, E denotes a complex Hilbert space.

DEFINITION 8. --- Let D be a vector subspace of E. A partial Hermitian form on E with domain D is a correspondence \(q = (\Gamma, \mathrm{E} \times \mathrm{E}, \mathbf{C})\) whose domain of definition is \(\mathrm{D} \times \mathrm{D}\) , whose graph \(\Gamma\) is functional, and such that the map from \(\mathrm{D} \times \mathrm{D}\) into \(\mathbf{C}\) defined by \(\Gamma\) is a Hermitian form. We denote \(\operatorname{dom}(q) = \mathrm{D}\) .

We say that \(q\) is a positive partial form if the Hermitian form it defines is positive.

Let u be a partial self-adjoint operator on E. For all elements x and y of E, we denote \(\mu_{x,y}\) (resp. \(\mu_{x}\) ) as the spectral measure of \((x,y)\) (resp. of x) relative to u.

Let \((j, |u|)\) be the polar decomposition of \(u\) (definition 7 of IV, p. 290). Let D' denote the domain of \(|u|^{1/2}\) . By definition, it is the space of \(x \in E\) such that the function \(z \mapsto |z|\) is integrable on \(\mathrm{Sp}(u)\) with respect to the measure \(\mu_x\) . It contains \(\mathrm{dom}(u)\) .

Let \((x,y)\in\mathrm{D}^{\prime}\times\mathrm{D}^{\prime}\) . The identity function of \(\mathrm{Sp}(u)\) is integrable with respect to the measure \(\mu_{x,y}\) (prop. 4, b) of IV, p. 271); we set

\[
q _ {u} (x, y) = \int_ {\operatorname{Sp} (u)} t d \mu_ {x, y} (t).
\]

If \(y \in \operatorname{dom}(u)\) , we have \(q_u(x, y) = \langle x \mid u(y) \rangle\) by formula (6) of IV, p. 271. The map \(q_u\) is a Hermitian form on D'. It defines a partial Hermitian form on E, said to be associated with u. If the operator u is positive, then the form \(q_u\) is a positive partial Hermitian form.

DEFINITION 9. --- Let \(q\) be a positive partial Hermitian form on E. We denote by \(\mathrm{E}_q\) the separated pre-Hilbert space \(\operatorname{dom}(q)\) equipped with the scalar product

\[
(x \mid y) _ {q} = \langle x \mid y \rangle + q (x, y).
\]

We denote by \(\|x\|_{q}\) the norm of \(x \in E_{q}\) . The form q is said to be closed if \(E_{q}\) is a Hilbert space.

PROPOSITION 13. --- Let u be a positive self-adjoint partial operator on E. Let \(q_{u}\) be the positive partial form associated with u.

\begin{enumerate}
\def\labelenumi{\alph{enumi})}
\item
  The domain of \(q_{u}\) is the domain of \(u^{1/2}\) , and for all x and y in \(\operatorname{dom}(u^{1/2})\) , we have \(q_{u}(x,y)=\langle u^{1/2}(x)\mid u^{1/2}(y)\rangle\) ;
\item
  The positive partial form \(q_{u}\) is closed. The domain of u is a core for \(q_{u}\) .
\end{enumerate}

Since \(u\) is positive, we have \(|u| = u\) . The domain \(\mathrm{dom}(|u|^{1/2})\) of \(q_u\) therefore coincides with \(\mathrm{dom}(u^{1/2})\) and we have

\[
\begin{array}{c} q _ {u} (x, y) = \int_ {\operatorname{Sp} (u)} t d \mu_ {x, y} (t) = \int_ {\operatorname{Sp} (u)} t ^ {1 / 2} \cdot t ^ {1 / 2} d \mu_ {x, y} (t) \\ = \langle u ^ {1 / 2} (x) | u ^ {1 / 2} (y) \rangle \end{array}
\]

for \(x\) and \(y\) in \(\mathrm{dom}(q_u)\) (prop. 4, b) of IV, p. 271). This proves the assertion \(a\) ).

Moreover, the pre-Hilbert space \(E_{q_{u}}\) then coincides with the pre-Hilbert space \(E_{u^{1/2}}\) associated with \(u^{1/2}\) (def. 4 of IV, p. 230). Since \(u^{1/2}\) is a closed partial operator, this space is a Hilbert space (prop. 4 of IV, p. 230) and \(\text{dom}(u)\) is dense in \(E_{u^{1/2}}\) by assertion \(e\) ) of Cor. 1 of IV, p. 273.

Let \(q\) be a partial Hermitian form on \(E\) with domain \(D\) dense in \(E\) . For \(y \in D\) , denote by \(\lambda_y\) the linear form \(x \mapsto q(y, x)\) on \(D\) . We denote by \(\widetilde{D}\) the set of \(y \in D\) such that the linear form \(\lambda_y\) is continuous on \(D\) .

Let \(y \in \widetilde{D}\) . Since D is dense in E, there exists a unique continuous linear form on E extending \(\lambda_{y}\) . We still denote it by \(\lambda_{y}\) . There exists a unique element \(u(y)\) in E such that \(\lambda_{y}(x) = \langle u(y) | x \rangle\) for all \(x \in E\) (EVT, V, p. 15, th. 3). The map \(y \mapsto u(y)\) is linear from \(\widetilde{D}\) into E; we say that the partial operator u on E with domain \(\widetilde{D}\) is the partial operator representing q. We therefore have

\[
q (x, y) = \langle x \mid u (y) \rangle
\]

for \(y \in \operatorname{dom}(u)\) and \(x \in D\) .

Remark. --- Let \(q\) be a closed positive partial form and \(q'\) a continuous positive Hermitian form on E. The positive Hermitian form defined on \(\text{dom}(q)\) by \((x, y) \mapsto q(x, y) + q'(x, y)\) is a closed positive partial Hermitian form with the same domain as \(q\) . We denote it by \(q + q'\) .

According to EVT, V, p. 16, cor. 2, there exists a unique linear map \(u' \in \mathcal{L}(\mathrm{E})\) such that \(q'(x,y) = \langle x \mid u'(y) \rangle\) for all \((x,y) \in \mathrm{E} \times \mathrm{E}\) . The endomorphism \(u'\) is positive and the partial operator representing the closed positive partial Hermitian form \(q + q'\) is \(u + u'\) .

PROPOSITION 14. --- Let q be a closed positive partial form with dense domain on E and let u be the partial operator representing q.

\begin{enumerate}
\def\labelenumi{\alph{enumi})}
\item
  The domain of u is dense in the Hilbert space \(E_{q}\) ;
\item
  The partial operator u is self-adjoint and positive.
\end{enumerate}

Since q is closed, the pre-Hilbert space \(E_{q}\) of Definition 9 is a Hilbert space.

Let us show that the partial operator \(u + 1_{E}\) with domain \(\operatorname{dom}(u)\) is bijective. Since

\[
\langle x \mid (u + 1 _ {\mathrm{E}}) (x) \rangle = q (x, x) + \| x \| ^ {2} \geqslant \| x \| ^ {2}
\]

for all \(x \in \operatorname{dom}(u)\) , this partial operator is injective.

Let \(y \in E\) . The linear form on \(E_{q}\) defined by \(x \mapsto \langle y | x \rangle\) is continuous since \(\|x\| \leqslant \|x\|_{q}\) . There therefore exists \(y' \in E_{q}\) such that

\[
\langle y \mid x \rangle = (y ^ {\prime} \mid x) _ {q} = \langle y ^ {\prime} \mid x \rangle + q (y ^ {\prime}, x)
\]

for all \(x \in E_{q}\) (EVT, V, p. 15, th. 3). By definition, this means that \(y'\) belongs to the domain of the partial operator \(\widetilde{u}\) representing the partial form \((x, y) \mapsto (x \mid y)_{q}\) with domain dom(q), and that \(\widetilde{u}(y') = y\) . Since \(\widetilde{u} = u + 1_{E}\) by the remark above, we deduce from this that the partial operator \(u + 1_{E}\) is also surjective, hence bijective.

Let us prove that the domain of \(u\) is dense in \(\mathrm{E}_q\) . Let \(y \in \mathrm{E}_q\) be orthogonal to \(\operatorname{dom}(u)\) in \(\mathrm{E}_q\) . There exists \(y' \in \operatorname{dom}(u)\) such that \(u(y') + y' = y\) . We then have

\[
0 = (y \mid y ^ {\prime}) _ {q} = \langle y \mid y ^ {\prime} \rangle + q (y, y ^ {\prime}) = \langle y \mid y ^ {\prime} + u (y ^ {\prime}) \rangle = \| y \| ^ {2}
\]

hence \(y = 0\) . The assertion \(a\) ) is therefore proved.

Since \(\mathrm{dom}(q)\) is dense in E and \(\| x\| \leqslant \| x\| _q\) for all \(x\in \mathrm{dom}(q)\) , assertion \(a\) ) implies that \(\mathrm{dom}(u)\) is dense in E.

Since the form \(q\) is Hermitian (resp. positive), for all \(x\) and \(y\) in \(\operatorname{dom}(u)\) , we have

\[
\langle y \mid u (x) \rangle = q (y, x) = \overline {{q (x , y)}} = \overline {{\langle x \mid u (y) \rangle}} = \langle u (y) \mid x \rangle ,
\]

(resp. \(\langle x \mid u(x) \rangle = q(x, x) \geqslant 0\) ) so that u is symmetric (resp. positive). Finally, the partial operator \(u + 1_{E}\) is self-adjoint by the corollary of Proposition 10 of IV, p. 240, and the same is true of u.

THEOREM 3. --- The map that associates the positive partial form \(q_u\) with a positive self-adjoint operator u on E is a bijection between the set of positive self-adjoint partial operators on E and the set of closed positive partial forms with dense domain on E. The inverse bijection associates with a positive partial form q the partial operator representing q.

By Propositions 13, b) and 14, b), the maps described in the statement are well defined. Let us show that they are inverse bijections of one another.

Let \(u\) be a positive self-adjoint partial operator on E and \(q\) the positive partial form associated with \(u\) . Denote by \(v\) the positive self-adjoint operator representing \(q\) . Let \(y \in \mathrm{dom}(u) \subset \mathrm{dom}(u^{1/2}) = \mathrm{dom}(q)\) . For every \(x \in \mathrm{dom}(q) = \mathrm{dom}(u^{1/2})\) , we have

\[
q (y, x) = \langle u ^ {1 / 2} (y) | u ^ {1 / 2} (x) \rangle = \langle u (y) | x \rangle ,
\]

hence y belongs to the domain of v and satisfies \(v(y) = u(y)\) . The partial operator v is therefore an extension of u; since u and v are self-adjoint, they are equal.

Conversely, let q be a closed positive partial form with dense domain and let u be the positive self-adjoint partial operator representing q. We have \(\operatorname{dom}(u) \subset \operatorname{dom}(u^{1/2})\) . For x and y in \(\operatorname{dom}(u)\) , we have

\[
q _ {u} (x, y) = \langle u ^ {1 / 2} (x) | u ^ {1 / 2} (y) \rangle = \langle x | u (y) \rangle = q (x, y)
\]

by Proposition 13, a). Thus, the Hilbert spaces \(\mathrm{E}_q\) and \(\mathrm{E}_{q_u}\) both contain \(\operatorname{dom}(u)\) as a dense subspace (Prop. 14, a) and Prop. 13, b), respectively) and \(\|x\|_q = \|x\|_{q_u}\) for all \(x \in \operatorname{dom}(u)\) . It follows that \(\mathrm{E}_q = \mathrm{E}_{q_u}\) and that \(q = q_u\) (lemma 8).

COROLLARY. --- Let \(q\) be a closed positive partial form on \(E\) and \(u\) the positive self-adjoint operator representing \(q\) . The domain of \(q\) is equal to the domain of \((1_{E} + u)^{1/2}\) , and we have

\[
\| x \| _ {q} = \| (1 _ {\mathrm{E}} + u) ^ {1 / 2} x \|
\]

for all \(x \in \operatorname{dom}(q)\) .

The domain of \(q\) is equal to the domain of \(u^{1/2}\) (prop. 13, \(a\) ), which coincides with that of \((1_{\mathrm{E}} + u)^{1/2}\) (cf. prop. 3 of IV, p. 271). For every \(x \in \operatorname{dom}(u) \subset \operatorname{dom}(u^{1/2})\) , we have

\[
\| (1 _ {\mathrm{E}} + u) ^ {1 / 2} x \| ^ {2} = \langle x | (1 _ {\mathrm{E}} + u) (x) \rangle = \| x \| ^ {2} + \langle x | u (x) \rangle = \| x \| _ {q} ^ {2}.
\]

As the domain of \(u\) is dense in the Hilbert space \(\mathrm{E}_q\) (prop. 14, a)), this formula extends by continuity to all \(x \in \mathrm{dom}(u^{1/2})\) .

Example. --- Let u be a positive partial operator on E which is not necessarily closed. One defines a positive partial form q with domain dom(u) by

\[
q (x, y) = \langle x \mid u (y) \rangle
\]

for \(x\) and \(y\) in \(\mathrm{dom}(u)\) . Let \(\mathrm{E}_q\) be the separated pre-Hilbert space of Definition 9 of IV, p. 292. Let \(\widetilde{\mathrm{E}}_q\) be the completed Hilbert space of \(\mathrm{E}_q\) , still denoted \((x,y) \mapsto (x \mid y)_q\) for the scalar product. Since the canonical inclusion \(\iota\) of \(\mathrm{E}_q\) in E is continuous, it admits a unique continuous extension, denoted \(\widetilde{\iota}\) , of \(\widetilde{\mathrm{E}}_q\) in E. The Hermitian form \(q\) is continuous on \(\mathrm{E}_q\) , hence also extends to a unique continuous positive Hermitian form \(\widetilde{q}\) on \(\widetilde{\mathrm{E}}_q\) .

Let us prove that the linear map \(\widetilde{\iota}\) is injective. Let \(x\in\mathrm{Ker}(\widetilde{\iota})\) . Consider a sequence \((x_{n})_{n\in\mathbf{N}}\) in \(\mathrm{E}_{q}\) which converges to \(x\) in \(\widetilde{\mathrm{E}}_{q}\) . We then have

\[
\lim _ {n \to + \infty} \iota (x _ {n}) = \widetilde {\iota} (x) = 0,
\]

therefore the sequence \((x_{n})_{n\in\mathbf{N}}\) converges to 0 in E. Let \(y\in E_{q}\) . Since \(\widetilde{q}\) is continuous on \(\widetilde{E}_{q}\) , it follows that

\[
\begin{array}{r l} (x \mid y) _ {q} = \lim _ {n \to + \infty} (x _ {n} \mid y) _ {q} = & \lim _ {n \to + \infty} \Big (\langle x _ {n} \mid y \rangle + q (x _ {n}, y) \Big) \\ & = \lim _ {n \to + \infty} \Big (\langle x _ {n} \mid y \rangle + \langle x _ {n} \mid u (y) \rangle \Big) = 0. \end{array}
\]

As the space \(E_{q}\) is dense in \(\widetilde{E}_{q}\) , we deduce that x = 0, as desired.

Identifying \(\widetilde{E}_{q}\) with its image under \(\widetilde{\iota}\) in E, we interpret \(\widetilde{q}\) as a closed positive partial form with dense domain extending q. The positive self-adjoint operator associated with \(\widetilde{q}\) is a self-adjoint extension of u. It is called the Friedrichs extension of u.

Note. --- Let \(c \in \mathbf R_{+}\) . Let q be a closed, densely defined partial Hermitian form such that \(q(x, x) \geqslant -c \|x\|^2\) for all x belonging to the domain of q. This means that the closed partial Hermitian form with domain dom(q) defined by \(\widetilde{q}(x, y) = q(x, y) + c \langle x \mid y \rangle\) is a positive partial form. By Theorem 3, such forms therefore correspond to self-adjoint partial operators u on E such that \(u + c1_{E}\) is positive, that is, such that u is bounded below (def. 7 of IV, p. 237) by -c.

Conversely, let u be a symmetric partial operator on E such that

\[
\langle x \mid u (x) \rangle \geqslant - c \| x \| ^ {2}\tag{17}
\]

for all \(x \in \operatorname{dom}(u)\) . Then \(v = u + c1_{E}\) is a positive partial operator on E. The Friedrichs extension of u is the self-adjoint operator \(\widetilde{v} - c1_{E}\) , where \(\widetilde{v}\) is the Friedrichs extension of v; it is a self-adjoint extension of u that does not depend on the choice of the real number c satisfying (17).
% semantic-fragments: legacy-input-seam
\relax{}
